% ----------------------------------
%   Appendix B: Local Volatility Surface — Implementation Details
% ----------------------------------

\chapter{Local Volatility Surface - Implementation Details}
\label{app:local_vol}


This appendix collects the implementation of the local volatility calibration of
Chapter~\ref{cha:cal_local_vol} and the decisions taken along the way that the main text
states without justifying.
All code lives in \href{https://github.com/M315/TFM/tree/main/scripts/local_vol_egger/}{Local Volatility Code}
and runs from the same cached SPY snapshot dated 2026-04-30
used throughout Appendix~\ref{app:impl_vol}.
The first sections describe the discrete problem that is actually solved, the domain, the
parametrization of the surface, the observation operator and the gradient sweep.
The later ones record the choices that could reasonably have been made differently, the prior,
the term structure, the selection of expirations, the regularization parameters and the
normalization of the unknown.

% ----------------------------------
\section{Discretization and computational domain}
\label{subsec:discretization_details}

The discretization is the one announced in Section~\ref{sec:variational_forms}, $P_1$ elements
in $y$ and backward Euler in $\tau$. What the main text leaves open is the size of everything.

The spatial domain is not fixed in advance, it is read off the data. Each kept expiration
contributes the log-moneyness range of its quoted strikes inside the band
$0.8 \le K/S_0 \le 1.2$, the union of those ranges is $y \in [-0.217,\,0.177]$, and the
computational domain is that union padded by $0.25$ on each side,
$y \in [-0.467,\,0.427]$, that is strikes from $0.63\,S_0$ to $1.53\,S_0$.
The pad is what separates the artificial boundary from the data.
Its width is about $1.8$ standard deviations of the longest maturity,
$\sigma\sqrt{\tau} \approx 0.14$ at $\sigma = 0.15$ and $\tau = 0.88$,
so whatever error we commit at the boundary is damped by the diffusion before it reaches the
observed band. Widening it further only buys solve time, since the coefficient out there is
constrained by no quote at all.

The mesh is uniform, $199$ elements and therefore $200$ nodes, a spacing of
$h \approx 4.5\times 10^{-3}$ in $y$, under half a percent in strike and finer than the quoted
strike spacing everywhere in the band. Time is divided into $300$ steps up to the last
maturity, $\Delta\tau \approx 2.9\times 10^{-3}$ years, close to one calendar day.

Backward Euler is a deliberate choice over Crank-Nicolson. Crank-Nicolson is second order in
time, but it is not strongly damping: its amplification factor tends to $-1$ for the stiff
modes, so a non-smooth initial condition leaves oscillations that decay only slowly. Our
initial condition is the call payoff, whose kink at $y = 0$ excites exactly those modes, and
the oscillations are worst in the second derivative in strike, which is the quantity carrying
the information about $a$. The standard cure is a Rannacher startup of a few implicit steps
\cite{giles_carter_crank_nicolson}. We use implicit Euler throughout instead, which is
uniformly damping at first order, because the time error is not the binding constraint here:
Table~\ref{tab:discounted_pde} below measures it at a day-sized step as
$5\times 10^{-2}$ dollars on a spot of $720$, a relative $10^{-4}$, an order of magnitude
below the reprice errors we report.

The solver works in dollars, not in the normalized price $u = C/S_0$ of
Section~\ref{sec:change_variable}. The equation is linear so the two differ by a constant
factor and describe the same coefficient, but the data term scales as $S_0^2$ while
$\beta\|a - a^*\|^2_{H^1}$ does not, so $\beta$ is tied to the scale of the prices. The values
quoted below are for prices in dollars at a spot of $720$, and they cannot be carried to
another underlying, or to the normalized form, without rescaling.

The boundary conditions are Dirichlet at both edges. The limits recorded
in~\eqref{eq:inverse_pde_2}, intrinsic value on the left and zero on the right, hold as
$M \to \infty$ and are too crude at a pad of $0.25$, where a call struck at $0.63\,S_0$ is
still worth appreciably more than its intrinsic value at long maturity. We instead evaluate
the Black-Scholes price at the two edge strikes, using that maturity's zero rate and dividend
yield and a wing volatility interpolated across the term structure from the extreme quoted
strike of each smile. This is exact whenever the smile is flat outside the quoted band, which
is the same assumption already built into the flat extension of the prior. At $\tau = 0$ the
boundary value is the payoff, so the data are continuous at the corner.

% ----------------------------------
\section{The surface as a stack of time slices}
\label{subsec:slice_parametrization}

The unknown is the stack of spatial slices of~\eqref{eq:coefficient_slices}, one per quoted
maturity. With eleven maturities and $200$ nodes that is $2\,200$ numbers, held as a single
flat vector, which is what the optimizer sees and what the bound constraints act on.

The slice boundaries have to live on the time grid, so each maturity is snapped to the nearest
time level, and if two maturities land on the same level the later one is pushed forward by
one step. On this data the eleven maturities map to distinct levels and the longest lands
exactly on the final step. The maturity the model prices is therefore displaced from the
quoted one by at most half a step, about half a day, and we do not correct for it: the
displacement is well inside the uncertainty of an end of day quote.

Slice $k$ then covers $(\tau_{k-1},\tau_k]$ and the observation at $\tau_k$ is read off the
trajectory at exactly the level where slice $k$ ends. This is what makes the triangular
structure of Section~\ref{sec:time_dependent_a} exact in the discrete problem: the first
expiration sees only the first slice, the second one the first two, and so on.

The optimizer imposes $a \ge 10^{-8}$ and no upper bound, although the admissible
set~\eqref{eq:admissible_set} carries one. The lower bound is the one that matters, it keeps
the operator parabolic and the recovered volatility real, and it is what makes the calibrated
surface arbitrage-free by construction.

% ----------------------------------
\section{The residual, and how a smile becomes data}
\label{subsec:observation_operator}

The misfit is the finite element $L^2$ norm over the observed band, not a sum of squared
errors over the quotes. Each expiration's call prices are interpolated linearly onto the mesh
nodes, and the difference against the computed price is multiplied by the $\{0,1\}$ indicator
of the nodes lying inside that expiration's quoted range. Because the mask is a projection,
$P^2 = P$, the same masked residual is simultaneously the integrand of the data term and the
source injected into the adjoint sweep, and no separate transpose has to be worked out.

The consequence of measuring the misfit in a continuum norm is that quotes are weighted by
mesh spacing rather than by how many of them there are. The expiration with $116$ strikes in
the band and the one with $46$ contribute the same integrated weight over the same interval of
moneyness, and a region where the chain quotes every dollar does not outweigh one where it
quotes every five. A plain sum over quotes would tilt the fit toward the densely quoted
expirations and toward the money, for no modelling reason.

The mask is a hard cut at the ends of each quoted range. Nothing outside is fitted, so the
coefficient in the wings is set by the penalty alone, which is exactly the behaviour visible
in the recovered slices of Figure~\ref{fig:egger_slices}.

% ----------------------------------
\section{The gradient in a single backward sweep}
\label{subsec:gradient_sweep}

One evaluation of the functional and its gradient is the following sequence.

\begin{enumerate}[label=(\roman*)]
    \item March the forward equation from $\tau = 0$ to the last maturity, storing the whole
          trajectory. Here that is $301$ states of $200$ values, so everything fits in memory
          and no checkpointing is needed.
    \item At each observed level, form the masked residual, accumulate its squared norm into
          the data term, and set the residual aside as an injection.
    \item March the adjoint equation backward from the last level, started from the last
          maturity's residual.
    \item At each step, after solving for the new adjoint state, assemble that step's
          contribution to the gradient and add it to the slice that owns the step.
    \item Only then, add the residual observed at that level to the adjoint state.
    \item Add the gradient of the penalty.
\end{enumerate}

The order of (iv) and (v) is the part that is easy to get wrong. Injecting the residual of
$\tau_j$ before taking the step's gradient contribution would let the price quoted at $\tau_j$
feed the coefficient on $(\tau_j,\tau_{j+1}]$, that is, it would make an option depend on the
volatility after its own expiry. Taking the contribution first keeps the dependence causal and
reproduces the triangular structure at the discrete level.

Two further points deserve recording.

The gradient is assembled directly as a linear form against the test functions and stopped
there, with no mass matrix solve. Solving $M g = b$ would instead return the Riesz
representative of the derivative in $L^2$, which is the natural gradient of the continuous
problem, whereas the assembled vector $b$ is the derivative with respect to the nodal
coefficients. Both are legitimately called the gradient, in different inner products, and
mixing the two is the usual way to obtain a gradient test that passes on one mesh and fails on
the next. We use the coefficient version everywhere, and the penalty gradient is assembled to
match, as $2\beta(M + K)(a - a^*)$ with $M$ and $K$ the mass and stiffness matrices, so the
two terms of the functional are differentiated in the same sense. It is also the version
L-BFGS-B expects, since its line search and its box constraints act on the coefficients.

What is differentiated is the discrete forward operator, not the continuous one. Every term of
the variational form~\eqref{eq:variational_forward_eq} carrying $a$ contributes: the diffusion,
the $\partial_y a$ term produced by the integration by parts, and the $a\,u_y$ half of the
convection. This is why the duality test of Section~\ref{sec:adjoint_verification} closes to
rounding error, $2.3\times 10^{-15}$, rather than to discretization error. Discretizing the
continuous adjoint instead would have given a gradient that is correct in the limit and
slightly wrong on the grid, which is enough to stall a quasi-Newton method near the minimum.

The cost is two PDE solves per iteration whatever the number of slices and the number of
nodes. Calibrating the eleven slices of the surface costs the same per iteration as
calibrating the single curve of Example~1.

% ----------------------------------
\section{The discrete penalty}
\label{subsec:penalty_details}

The two penalties of~\eqref{eq:surface_functional} are assembled from the same matrices. The
$H^1$ term is the full norm, values and slopes, and no condition is imposed on $a$ at the edges
of the domain, so the minimizer meets the boundary with a flat coefficient, which is consistent
with the flat wings of the prior.

The penalty across maturities is an $L^2$ jump rather than a discrete derivative. The
coefficient is piecewise constant in $\tau$ by construction, so $\partial_\tau a$ does not
exist and $\|a_k - a_{k-1}\|^2_{L^2}$ is the natural discrete stand-in. It is deliberately not
divided by the gap $\tau_k - \tau_{k-1}$. The gaps in this chain run from twelve days to
seventy-nine, and dividing by them would let the surface move almost freely across the long
gaps, which is precisely where the data are thinnest and the smoothing is most needed.

Only the $H^1$ term sees the prior. The jump term is prior free, so it is a genuine smoothness
requirement in maturity rather than a second pull toward the seed.

% ----------------------------------
\section{The prior and the seed}
\label{subsec:prior_details}

Section~\ref{sec:market_surface} explains why the prior cannot be the implied variance
$\tfrac12\sigma_{\mathrm{iv}}^2$ and states the inversion~\eqref{eq:bbf_inversion} used
instead. It comes from the short maturity result of \cite{berestycki_busca_florent}, that the
implied volatility is the harmonic mean of the local volatility between spot and strike,
\begin{equation}\label{eq:bbf_harmonic}
    \frac{y}{\sigma_{\mathrm{iv}}(y)}
        \;=\; \int_0^y \frac{dx}{\sigma_{\mathrm{loc}}(x)}
    \qquad \text{as } \tau \to 0,
\end{equation}
whose derivative in $y$ gives~\eqref{eq:bbf_inversion} directly.

Two details of the inversion matter in practice. The smile is replaced by a quadratic fit over
the band before it is inverted, because the inversion needs
$\partial_y\sigma_{\mathrm{iv}}$ and differentiating a raw quoted smile amplifies quote noise
into the prior; a quadratic is the lowest order carrying both a skew and a curvature, so the
seed keeps the two features that matter without importing the wiggles of one day's chain. And
the denominator of~\eqref{eq:bbf_inversion} is floored at $0.35$. On the downside wing the
skew is steep enough for $1 - y\,\partial_y\sigma_{\mathrm{iv}}/\sigma_{\mathrm{iv}}$ to
collapse toward zero, and a short maturity asymptotic used at $\tau$ close to a year has no
business producing an unbounded seed; the floor caps it at roughly three times the implied
level. Outside the quoted band the seed is extended flat.

The same field is used as the prior $a^*$ and as the starting iterate, so the optimizer begins
at the centre of the penalty and moves away from it only where the prices ask it to.

% ----------------------------------
\section{Interest and dividend as piecewise constant forwards}
\label{subsec:term_structure_details}

The implied volatility pipeline delivers one zero rate $r_j$ and one dividend yield $q_j$ per
expiration. The forward equation needs an instantaneous rate at every time level, so both are
bootstrapped into forwards that are constant on each interval between quoted maturities and
chosen so that integrating up to $\tau_j$ returns $z_j\tau_j$. Each maturity is then priced
with exactly the discount factor that was used to invert its smile, and the pair $r - q$ is
Egger's time dependent drift. The boundary values are the one place where the zero rates are
used rather than the forwards, since those are Black-Scholes prices for a maturity and not
local coefficients.

Bootstrapping differentiates the curve, so it amplifies whatever noise the zero rates carry.
For $r$ this is harmless, the forwards stay between $3.5\%$ and $3.9\%$. For $q$ it is
visible: the per-expiration put-call parity estimates are noisy at the tenth of a percent
level, and the implied forwards swing between $-0.4\%$ and $3.6\%$, one interval slightly
negative. We leave them as the bootstrap produces them. Smoothing the dividend term structure
would be defensible, but the alternative of a single constant $q$ misprices every maturity's
forward by construction, and the per-interval wobble enters only through the drift $r - q$
over one interval at a time.

% ----------------------------------
\section{From the option chain to the smiles}
\label{subsec:market_data_details}

Everything upstream of the calibration is the pipeline of Appendix~\ref{app:impl_vol}, the same
snapshot, the same Treasury rates, the same put-call parity dividend yields and the same
inverted volatilities.

The prices handed to the calibration are call prices, since that is what the Dupire forward
equation solves for, but they are not the quoted calls. Below spot the quoted put is converted
by put-call parity, $C = P + S_0e^{-q\tau} - Ke^{-r\tau}$, so every price comes from the liquid
out of the money leg while still describing a call. This is the out of the money composite of
Appendix~\ref{app:impl_vol} re-expressed in call prices.

The selection of expirations is the one place where the two pipelines deliberately part.
Appendix~\ref{app:impl_vol} keeps the most liquid expiration in each of ten log-spaced maturity
buckets, which balances a surface that is only going to be drawn. Here the slices are coupled
in time, so skipping a liquid expiration is not neutral: it leaves the calibration to cross a
gap the market does not have, with the jump penalty rather than quotes deciding what happens
inside it. The buckets also widen at the long end, exactly where the chain is already sparse.
We therefore keep every expiration whose composite carries at least twenty five strikes inside
the band, and drop the two that sit within two weeks of a kept one, 31 August and 31 December
2026. Two expirations a few days apart carry nearly the same information, and the second slice
would be fixed by the penalty rather than by its own quotes. Eleven expirations survive, with
maturities from $0.096$ to $0.882$ years and between $46$ and $116$ strikes each in the band.

% ----------------------------------
\section{Choice of the regularization parameters on market data}
\label{subsec:beta_market}

Section~\ref{sec:beta_choice} describes Egger's decreasing sequence of $\beta$ stopped by the
Morozov discrepancy principle, which is what the synthetic examples use, where $\delta$ is
known because the noise was added by hand.

On market data $\delta$ is not known. We estimated it from a relative price noise model, half a
percent of each quoted price over the observed band, which is the order of a bid-ask spread on
liquid SPY options, and ran the schedule against it. The stopping rule is not trustworthy at
that point. The estimate is a guess rather than a measurement, and more importantly the gap
between the quotes and the model is not noise in the sense the principle assumes: it also
contains the early exercise premium neglected in Appendix~\ref{app:impl_vol}, the fact that the
quotes across an expiration are not simultaneous, and the error in $r$ and $q$. A stopping
$\beta$ inherits all of it, and moves with an estimate we cannot pin down.

We therefore fix $\beta$ by hand for the market runs, $\beta = 10^{-4}$ for the single maturity
fit and $\beta_y = 10^{-1}$, $\beta_\tau = 1$ for the surface, chosen by inspecting the
trade-off directly: small enough that the reprice error settles near the level of the bid-ask
spread, large enough that the slices do not develop structure the quotes cannot support. As
noted above, these numbers are attached to prices in dollars at a spot of $720$ and to this
mesh, and they are not dimensionless.

% ----------------------------------
\section{Reporting the reprice error}
\label{subsec:reprice_metric}

The reprice errors of Figure~\ref{fig:egger_reprice} are relative errors on prices, which are
only meaningful where the price is not almost zero. We therefore report them over the liquid
strikes of each maturity, those whose market price is at least $0.5\%$ of spot, about $\$3.60$.
Without that floor the far out of the money calls, worth a few cents, would dominate the
maximum and say nothing about the quality of the fit. The model price at a quoted strike is
read back by interpolating the computed solution from the mesh nodes. The unweighted $L^2$
misfit over the whole band is reported separately, so the two measures are not confused: they
do not have to move together, and in Table~\ref{tab:discounted_calibration} they do not.

% ----------------------------------
\section{Decisions in the synthetic examples}
\label{subsec:synthetic_decisions}

Three choices in the reproduction of Examples~1 and~2 of \cite{egger2005tikhonov} are worth
recording, since none of them is stated in the paper.

\paragraph{The spot.} The text of the paper says $S = 100$, but its tabulated option values, a
call worth $439$ at $K = 600$, and the description of test D, an underlying worth $1000\$$,
only make sense at $S_0 = 1000$. The recovered $a$ is invariant under a rescaling of the
prices, but the Tikhonov balance is not, the data term scales as $S_0^2$ and the penalty does
not, so reproducing the paper's $\beta$ requires its $S_0$. We use $S_0 = 1000$.

\paragraph{The noise model.} The paper adds $0.1\%$ noise to the data and remarks that at
$S_0 = 1000$ this amounts to an uncertainty of about $\pm 1\$$, of the size of a typical
bid-ask spread. The two readings are not the same: $0.1\%$ of $S_0$ is a flat $\pm 1\$$ band on
every price, $0.1\%$ of each price is proportional to it. We use the relative reading. A flat
band is a far larger perturbation on the cheap out of the money calls, where the price itself
is of the order of a dollar and the whole signal is destroyed, and it did not reproduce the
paper's reconstructions; the relative reading does.

\paragraph{The inverse crime.} The paper generates its synthetic data on a wider domain with a
finer discretization and a different time integrator, so that the operator being inverted is
not the operator that produced the data. We follow the same idea, generating the data of
Example~2 on $[-5,5]$ with four times the spatial and temporal resolution and interpolating
onto the reconstruction grid, but keep implicit Euler on both sides rather than switching
schemes. The refinement alone already separates the two operators.

% ----------------------------------
\section{The discounted normalization $u = e^{\int_t^T q\,ds}\,C$}
\label{subsec:discounted_normalization}

Throughout Section~\ref{sec:market_surface} the unknown is the call price itself, in dollars
(Section~\ref{subsec:discretization_details}), and the equation keeps the explicit zero-order
term of~\eqref{eq:inverse_pde_1}.
\cite{egger2005tikhonov} instead absorb the dividend factor into the unknown by setting
\begin{equation*}
    \tilde{u}(y,\tau) = e^{\int_t^T q(s)\,ds}\,u(y,\tau),
\end{equation*}
which cancels the $-q\,u$ term and leaves the shorter equation
\begin{equation*}
    -\tilde{u}_\tau + a\,(\tilde{u}_{yy} - \tilde{u}_y) + (q - r)\,\tilde{u}_y = 0 .
\end{equation*}
A deterministic factor relates the two normalizations, so they describe the same continuous
problem and coincide exactly when $q = 0$, as in Examples~1 and~2.
On the market surface $q$ comes from put-call parity (Appendix~\ref{subsec:div_yield_cal})
and reaches about $1.1\%$, so the only thing that can separate them is the discretization.
Which one discretizes better is not obvious.
The discounted form treats the exponential decay in $\tau$ exactly instead of leaving it to the
backward Euler scheme, but it pays for that by rescaling the data and the boundary values,
which reweights the residual across maturities.
We implemented both and compared them.
Both remain in the code as a switch on the solver, and the two tables below are reproducible
from it.

The discounted form needs three changes: drop the zero-order coefficient from the variational
forms~\eqref{eq:variational_forward_eq} and~\eqref{eq:variational_adjoint_eq}, multiply the
observed prices and the Dirichlet data by $e^{\int_0^\tau q}$, and divide that factor out again
when the calibrated surface is repriced.
The two solvers do agree on the continuous problem.
Holding $a$ fixed at the implied volatility seed and converting the discounted solution back to
prices, the two forward solves match to $9\times 10^{-3}$ dollars on a spot of $720$.
That is a relative difference of $10^{-4}$, the $O(\Delta\tau)$ gap we expect.

Refining the time grid tells us which one is closer to the exact solution.
We solved forward at the production grid of $N = 300$ time steps and compared against a
reference computed on the same time levels with $N = 2400$.
Table~\ref{tab:discounted_pde} reports the error averaged over the eleven quoted maturities, and
repeats the experiment at artificially inflated dividend yields.

\begin{table}[h]
    \centering
    \begin{tabular}{c c c c}
        \toprule
        $q$ & explicit $-q\,u$ & discounted & ratio \\
        \midrule
        market ($\leq 1.6\%$) & $5.51\times10^{-2}$ & $5.38\times10^{-2}$ & $1.02$ \\
        $2\%$  & $5.25\times10^{-2}$ & $5.22\times10^{-2}$ & $1.01$ \\
        $5\%$  & $4.44\times10^{-2}$ & $4.37\times10^{-2}$ & $1.01$ \\
        $10\%$ & $3.37\times10^{-2}$ & $3.28\times10^{-2}$ & $1.03$ \\
        $20\%$ & $4.93\times10^{-2}$ & $4.68\times10^{-2}$ & $1.05$ \\
        \bottomrule
    \end{tabular}
    \caption{Time discretization error of the forward solve at $N = 300$ time steps, in dollars
             against an $N = 2400$ reference on the same time levels, averaged over the eleven
             quoted maturities. The first row uses the calibrated dividend term structure, the
             others a flat override. The discounted form never gains more than $5\%$, because
             the backward Euler error is dominated by the kink of the payoff and by the
             convection term rather than by the zero-order term.}
    \label{tab:discounted_pde}
\end{table}

Even at a dividend yield of $20\%$, an order of magnitude above anything SPY pays, the gain is
a few percent.
The same holds for the inverse problem.
Running the full surface calibration~\eqref{eq:surface_functional} both ways gives surfaces
that differ by up to one volatility point, which looks large until the two effects are
separated.
Table~\ref{tab:discounted_calibration} reprices each calibrated surface with both solvers.
Swapping the solver moves the median error by $0.003$ percentage points, swapping the surface
moves it by about $0.01$.
The gap therefore measures the path the optimizer took through a flat valley of the Tikhonov
functional, not the normalization.
Neither surface is better than the other either, since the one with the lower median error over
the liquid strikes carries the higher unweighted data misfit.

\begin{table}[h]
    \centering
    \begin{tabular}{l l c c}
        \toprule
        calibrated with & repriced with & median error & data misfit (\$$^2$) \\
        \midrule
        discounted        & discounted        & $0.116\%$ & $0.02881$ \\
        discounted        & explicit $-q\,u$  & $0.119\%$ & $0.02928$ \\
        explicit $-q\,u$  & discounted        & $0.131\%$ & $0.02776$ \\
        explicit $-q\,u$  & explicit $-q\,u$  & $0.128\%$ & $0.02817$ \\
        \bottomrule
    \end{tabular}
    \caption{Each calibrated surface of Section~\ref{sec:market_surface} repriced with both
             solvers ($\beta_y = 10^{-1}$, $\beta_\tau = 1$, eleven maturities). The median error
             is taken over the liquid strikes of each maturity, the data misfit is the
             unweighted term of~\eqref{eq:surface_functional}. The two rows of each block differ
             only by the solver and barely move, so what separates the blocks is the optimizer
             path.}
    \label{tab:discounted_calibration}
\end{table}

Neither approach outperforms the other on this data, so we keep the explicit zero-order form as
the default. It compares the solver output directly against the quoted prices, with no
rescaling in between.
The discounted form is the better choice on an underlying with a large carry, a high dividend
single stock or an FX option, where the foreign rate plays the role of $q$ and the factor
$e^{\int_t^T q}$ is no longer close to one.
